\scanpage{104}
\mspage{97}
% archive sheet number: 103
le \unsure{``il suffit''|'il suffit'} dans b), qui \deleted{prouve} \inserted{assure} que $X=\unsure{\underline{\lim}|\varinjlim}X_i\in M$, qed.

\uline{NB} On a prouvé $b\Rightarrow a$, sans utiliser \unsure{l'hyp. faite|l'hyp.thèse|l'hypothèse} que tt objet de $\mathrm{Filt}^{\pi}$ est $\pi$-accessible. Mais on a utilisé cette hypothèse pour prouver que b \deleted{implique} (ou a) implique que $M$ est stable par $\underline{\lim}$ grandes devant $\pi$ (lemme 2, p. 90).

$a\overset{?}{\Longrightarrow}b$. Le ``il suffit'' de b) est conséquence immédiate de a). Je veux prouver que si, dans \deleted{le \unsure{même système}} \inserted{système \inserted{inductif}} \deleted{\unsure{parf.}} \inserted{$\pi$-adapté} $(X_i)_{i\in I}$ de limite $X$, on a $X\in M$, alors $I_M$ est cofinal dans $I$. Mais par hyp. a), on a encore
\[
\deleted{\unsure{A|M}}X=\varinjlim_JY_j,
\]
avec $J$ grand devant $\pi$, les $Y_j$ dans $M^{\pi}$. \deleted{Soit $i\in I$}, Admettons que $\mathrm{Filt}^{\pi}\subset E_\pi$ i.e tt élément de $\mathrm{Filt}^{\pi}$, en particulier les $X_i$, sont \deleted{bien} $\pi$-accessibles.

Je pense que je n'arriverai pas à prouver $a_\pi\Rightarrow b_\pi$, mais seulement ceci: Si $\pi'$ est un cardinal $>\pi$ (disons le successeur de $\pi$ \unsure{pris|$\xi$} dans l'\deleted{\unsure{ens}} \inserted{suite ordinale} des cardinaux) alors
\[
a_\pi\Longrightarrow b_{\pi'}.
\]
\deleted{Donc dans les arguments qui précèdent, on supposera $I$ grand devant $\pi'$, non seulement grand devant $\pi$. Soit \inserted{donc} $\alpha_0$ le plus petit ordinal qui soit grand devant $\pi$, il est de cardinal $\leq\pi'$ (car lorsque $\pi'>\pi$, le plus petit ordinal $\unsure{\alpha'_0|\alpha'}$ de cardinal $\pi'$ est grand dev. \unsure{à} $\pi$ car $\pi<\pi'$, \unsure{donc} $\alpha'\geq\alpha'_0$, donc $\pi'\unsure{\leq|=}\operatorname{card}\alpha'\geq\operatorname{card}\unsure{\alpha_0|\pi_0}$.}
